\chapter{Harmonic Functions}

\section{Mean Value Properties}

Throughout, $\Omega\subset\mathbb R^n$ is a connected domain and $\omega_n$
denotes the surface measure of the unit sphere.

\begin{definition}\label{hanlin:def-mean-value-properties}\dcref{Definition 1.1}\group{harmonic-functions}
\source{han-lin-lecture-notes:page-0007}
\lean{HanLinLectureNotes.Ch01.HasMeanValueProperties}
\leanok
For $u\in C(\Omega)$, the spherical and ball mean value properties are,
respectively,
\[
 u(x)=\frac{1}{\omega_n r^{n-1}}\int_{\partial B_r(x)}u\,d\sigma,
 \qquad
 u(x)=\frac{n}{\omega_n r^n}\int_{B_r(x)}u\,dy,
\]
for every ball $B_r(x)\subset\Omega$.
\end{definition}

\begin{remark}\label{hanlin:remark-mean-value-equivalent}\dcref{Remark 1.2}\group{harmonic-functions}
\source{han-lin-lecture-notes:page-0007}
The two identities in Definition~\ref{hanlin:def-mean-value-properties} are
equivalent: integrate the spherical identity in $r$, or differentiate the ball
identity, using
\[
u(x)r^{n-1}=\frac1{\omega_n}\int_{\partial B_r(x)}u\,d\sigma,
\qquad
u(x)r^n=\frac n{\omega_n}\int_{B_r(x)}u\,dy.
\]
\end{remark}

\begin{remark}\label{hanlin:remark-mean-value-equivalent-forms}\dcref{Remark 1.3}\group{harmonic-functions}
\source{han-lin-lecture-notes:page-0007}
After scaling to the unit ball, the same formulas read
\[
u(x)=\frac1{\omega_n}\int_{|\omega|=1}u(x+r\omega)\,d\sigma_\omega,
\qquad
u(x)=\frac n{\omega_n}\int_{|y|\le1}u(x+ry)\,dy.
\]
\end{remark}

\begin{theorem}\label{hanlin:thm-mean-value-max}\dcref{Theorem 1.4}\group{harmonic-functions}
\source{han-lin-lecture-notes:page-0008}
\uses{hanlin:def-mean-value-properties}
\lean{HanLinLectureNotes.Ch01.meanValue_extrema_mem_frontier_of_not_constant}
\leanok
If $u\in C(\overline\Omega)$ has the mean value property in $\Omega$, then a
maximum or minimum in the interior forces $u$ to be constant; otherwise every
maximum or minimum of $u$ attained on $\overline\Omega$ lies on $\partial\Omega$.
\end{theorem}
\begin{proof}
\sketch For an interior maximum $M$, the set $\Sigma=\{u=M\}$ is relatively
closed. At $x_0\in\Sigma$, the ball mean value identity and $u\le M$ imply
equality in an average bound, hence $u=M$ on every ball compactly contained in
$\Omega$. Thus $\Sigma$ is open, and connectedness gives $\Sigma=\Omega$.
Apply the same argument to $-u$ for minima.
\end{proof}

\begin{definition}\label{hanlin:def-harmonic}\dcref{Definition 1.5}\group{harmonic-functions}
\source{han-lin-lecture-notes:page-0008}
\lean{HanLinLectureNotes.Ch01.IsHarmonicOn}
\leanok
A function $u\in C^2(\Omega)$ is harmonic when $\Delta u=0$ in $\Omega$.
\end{definition}

\begin{theorem}\label{hanlin:thm-harmonic-mean-value}\dcref{Theorem 1.6}\group{harmonic-functions}
\source{han-lin-lecture-notes:page-0008}
\uses{hanlin:def-mean-value-properties, hanlin:def-harmonic}
\lean{HanLinLectureNotes.Ch01.IsHarmonicOn.hasMeanValueProperties}
\leanok
Every harmonic $u\in C^2(\Omega)$ satisfies the mean value properties.
\end{theorem}
\begin{proof}
\sketch Apply the divergence theorem on $B_\rho(x)$ to obtain
\[
\int_{B_\rho(x)}\Delta u=\rho^{n-1}\frac d{d\rho}
\int_{|\omega|=1}u(x+\rho\omega)\,d\sigma_\omega.
\]
The left side vanishes, so the spherical integral is independent of $\rho$.
Its limit at $\rho=0$ is $\omega_nu(x)$, giving the spherical identity; the
ball identity follows from Remark~\ref{hanlin:remark-mean-value-equivalent}.
\end{proof}

\begin{remark}\label{hanlin:remark-mvp-not-smooth}\dcref{Remark 1.7}\group{harmonic-functions}
\source{han-lin-lecture-notes:page-0008}
\uses{hanlin:thm-harmonic-mean-value}
The mean value property is stated for continuous functions, but it itself
implies the regularity needed for harmonicity.
\end{remark}

\begin{theorem}\label{hanlin:thm-mean-value-harmonic}\dcref{Theorem 1.8}\group{harmonic-functions}
\source{han-lin-lecture-notes:page-0008}
\uses{hanlin:def-mean-value-properties, hanlin:thm-harmonic-mean-value}
If $u\in C(\Omega)$ has the mean value property, then $u\in C^\infty(\Omega)$
and $\Delta u=0$.
\end{theorem}
\begin{proof}
\sketch Let $\varphi\in C_0^\infty(B_1)$ be radial with integral one and set
$\varphi_\varepsilon(z)=\varepsilon^{-n}\varphi(z/\varepsilon)$. For
$\varepsilon<\operatorname{dist}(x,\partial\Omega)$, polar coordinates and
the mean value identity give $(\varphi_\varepsilon*u)(x)=u(x)$. Hence $u$ is
smooth by convolution. Differentiating the spherical identity and applying the
divergence theorem yields $\int_{B_r(x)}\Delta u=0$ for every admissible ball,
so $\Delta u=0$.
\end{proof}

\begin{remark}\label{hanlin:remark-mvp-harmonic-smooth}\dcref{Remark 1.9}\group{harmonic-functions}
\source{han-lin-lecture-notes:page-0009}
\uses{hanlin:thm-mean-value-max, hanlin:thm-harmonic-mean-value, hanlin:thm-mean-value-harmonic}
The preceding results give smoothness, the mean value property, and the maximum
principle for harmonic functions. Thus the bounded-domain Dirichlet problem
has at most one solution. On unbounded domains uniqueness can fail: on
$\{|x|>1\}$ the zero-boundary harmonic functions
\[
u(x)=\log|x|\ (n=2),\qquad u(x)=|x|^{2-n}-1\ (n\ge3)
\]
are nonzero, and $u(x)=x_n$ is a nonzero solution on the upper half-space.
\end{remark}

\begin{lemma}\label{hanlin:lemma-interior-gradient-1}\dcref{Lemma 1.10}\group{harmonic-functions}
\source{han-lin-lecture-notes:page-0010}
\uses{hanlin:thm-harmonic-mean-value, hanlin:lemma-interior-gradient-2}
\lean{HanLinLectureNotes.Ch01.IsHarmonicOn.fderiv_norm_le_ball}
\leanok
If $u\in C(\overline{B_R(x_0)})$ is harmonic in $B_R(x_0)$, then
\[
|Du(x_0)|\le\frac nR\max_{\overline{B_R(x_0)}}|u|.
\]
\end{lemma}
\begin{proof}
\sketch For any unit vector $e$, the directional derivative $D_eu$ is harmonic.
Its ball mean value formula and the divergence theorem give
\[
D_eu(x_0)=\frac n{\omega_nR^n}\int_{\partial B_R(x_0)}u\,(e\cdot\nu)\,d\sigma,
\]
and hence $|D_eu(x_0)|\le nR^{-1}\max|u|$. Take $e$ parallel to $Du(x_0)$.
\end{proof}

\begin{lemma}\label{hanlin:lemma-interior-gradient-2}\dcref{Lemma 1.11}\group{harmonic-functions}
\source{han-lin-lecture-notes:page-0010}
\uses{hanlin:thm-harmonic-mean-value}
\lean{HanLinLectureNotes.Ch01.IsHarmonicOn.fderiv_norm_le_of_nonnegative_ball}
\leanok
If $u\ge0$ is harmonic in $B_R(x_0)$ and continuous on its closure, then
\[
|Du(x_0)|\le\frac nR u(x_0).
\]
\end{lemma}
\begin{proof}
\sketch Use the boundary formula from Lemma~\ref{hanlin:lemma-interior-gradient-1},
replace $|\nu_i|$ by $1$, and use nonnegativity together with the mean value
identity for the boundary integral.
\end{proof}

\begin{corollary}\label{hanlin:corollary-liouville}\dcref{Corollary 1.12}\group{harmonic-functions}
\source{han-lin-lecture-notes:page-0010}
\uses{hanlin:lemma-interior-gradient-2}
\lean{HanLinLectureNotes.Ch01.IsHarmonicOn.eq_const_of_bddAbove_or_bddBelow}
\leanok
A harmonic function on $\mathbb R^n$ that is bounded above or bounded below is
constant.
\end{corollary}
\begin{proof}
\sketch Add a constant and change sign if necessary so that $u\ge0$. Apply
Lemma~\ref{hanlin:lemma-interior-gradient-2} on arbitrarily large balls centered
at a fixed point and let their radii tend to infinity; then $Du=0$ everywhere.
\end{proof}

\begin{lemma}\label{hanlin:lemma-higher-derivative-estimate}\dcref{Lemma 1.13}\group{harmonic-functions}
\source{han-lin-lecture-notes:page-0010}
\uses{hanlin:lemma-interior-gradient-1}
If $u$ is harmonic in $B_R(x_0)$ and continuous on the closure, then for every
multi-index $\alpha$ of order $m$,
\[
|D^\alpha u(x_0)|\le\frac{n^m e^{m-1}m!}{R^m}
\max_{\overline{B_R(x_0)}}|u|.
\]
\end{lemma}
\begin{proof}
\sketch Induct on $m$. Apply Lemma~\ref{hanlin:lemma-interior-gradient-1} to an
$m$th derivative on a concentric smaller ball, use the induction estimate on
the remaining radius, and choose $\theta=m/(m+1)$ in the radius split. The
factor $\theta^{-m}(1-\theta)^{-1}<e(m+1)$, and
$\alpha_1!\cdots\alpha_n!\le m!$ gives the stated multi-index bound.
\end{proof}

\begin{theorem}\label{hanlin:thm-harmonic-analytic}\dcref{Theorem 1.14}\group{harmonic-functions}
\source{han-lin-lecture-notes:page-0011}
\uses{hanlin:lemma-higher-derivative-estimate}
Every harmonic function is real analytic.
\end{theorem}
\begin{proof}
\sketch Taylor-expand in a ball $B_{2R}(x)$. Lemma~\ref{hanlin:lemma-higher-derivative-estimate}
bounds the order-$m$ remainder by a geometric factor
$((n^2e)|h|/R)^m\max_{B_{2R}(x)}|u|$, which tends to zero for sufficiently small
$|h|$. Thus the Taylor series converges to $u$ locally.
\end{proof}

\begin{theorem}\label{hanlin:thm-harnack}\dcref{Theorem 1.15}\group{harmonic-functions}
\source{han-lin-lecture-notes:page-0011}
\uses{hanlin:thm-harmonic-mean-value}
If $u\ge0$ is harmonic in $\Omega$, then for every compact $K\subset\Omega$ there
is $C=C(\Omega,K)$ such that
\[
C^{-1}u(y)\le u(x)\le Cu(y)\qquad(x,y\in K).
\]
\end{theorem}
\begin{proof}
\sketch The mean value formula compares values at any two points in one ball
$B_R$ whose fourfold enlargement lies in $\Omega$, with a factor depending
only on $n$. Connect points of $K$ by a finite chain of such balls and iterate
the local comparison.
\end{proof}

For later use, integration by parts gives
\[
\int_\Omega u\,\Delta\varphi=0\qquad(\varphi\in C_0^2(\Omega))
\]
for every harmonic $u$.

\begin{theorem}\label{hanlin:thm-distributional-harmonic}\dcref{Theorem 1.16}\group{harmonic-functions}
\source{han-lin-lecture-notes:page-0012}
\uses{hanlin:thm-mean-value-harmonic}
If $u\in C(\Omega)$ satisfies
\[
\int_\Omega u\,\Delta\varphi=0\qquad(\varphi\in C_0^2(\Omega)),
\]
then $u$ is harmonic in $\Omega$.
\end{theorem}
\begin{proof}
\sketch Test the distributional equation with compactly supported radial
functions in a ball. Differentiating the resulting identities in the radius
shows that the spherical mean of $u$ is constant in the radius; its limit at
zero is $u(x)$, so $u$ has the mean value property. Theorem~\ref{hanlin:thm-mean-value-harmonic}
then gives harmonicity. (For the explicit test-function induction, use powers
of $|y|^2-r^2$ and differentiate successively in $r$.)
\end{proof}

\section{Fundamental Solutions}

A radial harmonic function $u(x)=v(|x-a|)$ satisfies
$v''+(n-1)v'/r=0$. Hence $v=c_1+c_2\log r$ for $n=2$ and
$v=c_3+c_4r^{2-n}$ for $n\ge3$. Normalizing the outward flux to one gives
\[
\Gamma(a,x)=
\begin{cases}
 (2\pi)^{-1}\log|a-x|,&n=2,\\
 [\omega_n(2-n)]^{-1}|a-x|^{2-n},&n\ge3.
\end{cases}
\]
For $x\ne a$, $\Delta_x\Gamma(a,x)=0$, and
$\int_{\partial B_r(a)}\partial_{n_x}\Gamma\,d\sigma_x=1$.

\begin{theorem}\label{hanlin:thm-green-identity}\dcref{Theorem 1.17}\group{harmonic-functions}
\source{han-lin-lecture-notes:page-0014}
If $\Omega$ is bounded and $u\in C^1(\overline\Omega)\cap C^2(\Omega)$, then
for $a\in\Omega$
\[
u(a)=\int_\Omega\Gamma(a,x)\Delta u(x)\,dx-
\int_{\partial\Omega}\left(\Gamma(a,x)\partial_{n_x}u(x)-u(x)\partial_{n_x}\Gamma(a,x)\right)d\sigma_x.
\]
\end{theorem}
\begin{proof}
\sketch Apply Green's identity on $\Omega\setminus B_r(a)$ and let $r\downarrow0$.
The singular flux of $\Gamma$ contributes $u(a)$, while the other inner-boundary
term vanishes; the logarithmic case is analogous.
\end{proof}

\begin{remark}\label{hanlin:remark-green-identity}\dcref{Remark 1.18}\group{harmonic-functions}
\source{han-lin-lecture-notes:page-0014}
The singularity of $\Gamma(a,\cdot)$ is integrable. If $a\notin\overline\Omega$,
the right side of Theorem~\ref{hanlin:thm-green-identity} is zero. Setting
$u\equiv1$ gives
\[
\int_{\partial\Omega}\partial_{n_x}\Gamma(a,x)\,d\sigma_x=1\qquad(a\in\Omega).
\]
\end{remark}

\begin{remark}\label{hanlin:remark-local-green-gradient}\dcref{Remark 1.19}\group{harmonic-functions}
\source{han-lin-lecture-notes:page-0014}
\uses{hanlin:thm-green-identity}
Applying a cutoff version of Theorem~\ref{hanlin:thm-green-identity} to a
harmonic $u$ on $B_1$ yields the local estimates
\[
\sup_{B_{1/2}}|u|\le c\|u\|_{L^p(B_1)},\qquad
\sup_{B_{1/2}}|Du|\le c\|u\|_{L^\infty(B_1)},
\]
with $c$ depending only on $n$ (and $p$ in the first estimate).
\end{remark}

For the Dirichlet problem $\Delta u=f$ in $\Omega$, $u=\varphi$ on
$\partial\Omega$, add to $\Gamma$ the harmonic correction $\Phi(x,\cdot)$
with boundary data $-\Gamma(x,\cdot)$. The Green function
$G=\Gamma+\Phi$ then satisfies
\[
u(x)=\int_\Omega G(x,y)f(y)\,dy+
\int_{\partial\Omega}\varphi(y)\partial_{n_y}G(x,y)\,d\sigma_y,
\]
and is uniquely determined by the maximum principle.

\begin{lemma}\label{hanlin:lemma-green-symmetric}\dcref{Lemma 1.20}\group{harmonic-functions}
\source{han-lin-lecture-notes:page-0015,han-lin-lecture-notes:page-0016}
\uses{hanlin:thm-green-identity}
The Green function is symmetric: $G(x,y)=G(y,x)$ for distinct $x,y\in\Omega$.
\end{lemma}
\begin{proof}
\sketch Apply Green's identity to $G(x_1,\cdot)$ and $G(x_2,\cdot)$ on the
domain with small balls about both poles removed. The outer boundary term is
zero. Replacing each pole term by the normalized fundamental solution and
letting the radii tend to zero gives $G(x_1,x_2)=G(x_2,x_1)$.
\end{proof}

\begin{proposition}\label{hanlin:prop-green-sign}\dcref{Proposition 1.21}\group{harmonic-functions}
\source{han-lin-lecture-notes:page-0016}
\uses{hanlin:thm-mean-value-max}
For the Green function and $x\ne y$,
\[
0>G(x,y)>\Gamma(x,y)\quad(n\ge3),
\]
while for $n=2$,
\[
0>G(x,y)>\Gamma(x,y)-(2\pi)^{-1}\log\operatorname{diam}(\Omega).
\]
\end{proposition}
\begin{proof}
\sketch The maximum principle applied away from a small pole gives $G<0$.
Writing $G=\Gamma+\Phi$, the boundary values of the harmonic correction are
positive for $n\ge3$ and bounded below by $-(2\pi)^{-1}\log\operatorname{diam}(\Omega)$
for $n=2$; apply the maximum principle to $\Phi$.
\end{proof}

\begin{theorem}\label{hanlin:thm-green-ball}\dcref{Theorem 1.22}\group{harmonic-functions}
\source{han-lin-lecture-notes:page-0017}
For $x,y\in B_R$ the Green function is
\[
G(x,y)=\begin{cases}
[(2-n)\omega_n]^{-1}\left(|x-y|^{2-n}-\left|\frac R{|x|}x-\frac{|x|}{R}y\right|^{2-n}\right),&n\ge3,\\
(2\pi)^{-1}\left(\log|x-y|-\log\left|\frac R{|x|}x-\frac{|x|}{R}y\right|\right),&n=2.
\end{cases}
\]
\end{theorem}
\begin{proof}
\sketch Invert $x$ through $\partial B_R$ at $X=R^2x/|x|^2$. For $|y|=R$,
$|y-x|=(|x|/R)|y-X|$, so the displayed reflected singularity cancels the
boundary value while leaving the pole at $x$.
\end{proof}

\begin{corollary}\label{hanlin:cor-green-normal-derivative}\dcref{Corollary 1.23}\group{harmonic-functions}
\source{han-lin-lecture-notes:page-0017}
\uses{hanlin:thm-green-ball}
For $x\in B_R$ and $y\in\partial B_R$,
\[
\partial_{n_y}G(x,y)=\frac{R^2-|x|^2}{\omega_nR|x-y|^n}.
\]
\end{corollary}
\begin{proof}
\sketch Differentiate the formula in Theorem~\ref{hanlin:thm-green-ball} and use
$n_y=y/R$ together with the inversion identity on the boundary.
\end{proof}

Let $K(x,y)=\partial_{n_y}G(x,y)$ on $B_R\times\partial B_R$. The explicit
formula gives smoothness off the diagonal, positivity, vanishing near a fixed
boundary point away from that point, and $\Delta_xK=0$; integrating the Green
representation for $u\equiv1$ gives $\int_{\partial B_R}K(x,y)\,d\sigma_y=1$.

\begin{theorem}\label{hanlin:thm-poisson-dirichlet}\dcref{Theorem 1.24}\group{harmonic-functions}
\source{han-lin-lecture-notes:page-0018}
\uses{hanlin:cor-green-normal-derivative}
For $\varphi\in C(\partial B_R)$, define
\[
u(x)=\int_{\partial B_R}K(x,y)\varphi(y)\,d\sigma_y,\qquad |x|<R.
\]
Then $u$ is smooth and harmonic in $B_R$, continuous on $\overline{B_R}$, and
$u=\varphi$ on $\partial B_R$.
\end{theorem}
\begin{proof}
\sketch Smoothness and harmonicity follow from the kernel properties. For
boundary continuity, split the integral near and away from a boundary point;
continuity of $\varphi$ controls the near part and uniform kernel decay controls
the far part. The normalization $\int K=1$ completes the limit.
\end{proof}

\begin{remark}\label{hanlin:remark-poisson-mean-value}\dcref{Remark 1.25}\group{harmonic-functions}
\source{han-lin-lecture-notes:page-0019}
\uses{hanlin:thm-poisson-dirichlet}
At the center, the Poisson formula is the spherical mean value formula:
\[
u(0)=\frac1{\omega_nR^{n-1}}\int_{\partial B_R}u\,d\sigma.
\]
\end{remark}

\begin{lemma}\label{hanlin:lemma-harnack-ball}\dcref{Lemma 1.26}\group{harmonic-functions}
\source{han-lin-lecture-notes:page-0019}
\uses{hanlin:thm-poisson-dirichlet}
If $u\ge0$ is harmonic in $B_R(x_0)$ and $r=|x-x_0|<R$, then
\[
\left(\frac R{R+r}\right)^{n-2}\frac{R-r}{R+r}u(x_0)
\le u(x)\le
\left(\frac R{R-r}\right)^{n-2}\frac{R+r}{R-r}u(x_0).
\]
\end{lemma}
\begin{proof}
\sketch Translate to $x_0=0$, insert the Poisson kernel, and use
$R-r\le|y-x|\le R+r$ on $\partial B_R$. Replace the boundary integral by
$\omega_nR^{n-1}u(0)$ via the mean value formula.
\end{proof}

\begin{corollary}\label{hanlin:corollary-liouville-poisson}\dcref{Corollary 1.27}\group{harmonic-functions}
\source{han-lin-lecture-notes:page-0019}
\uses{hanlin:lemma-harnack-ball}
A harmonic function on $\mathbb R^n$ bounded above or below is constant.
\end{corollary}
\begin{proof}
\sketch Normalize to $u\ge0$ and apply Lemma~\ref{hanlin:lemma-harnack-ball} on
$B_R(0)$ for fixed $x$. Let $R\to\infty$; both comparison factors tend to one.
\end{proof}

\begin{theorem}\label{hanlin:thm-removable-singularity}\dcref{Theorem 1.28}\group{harmonic-functions}
\source{han-lin-lecture-notes:page-0019,han-lin-lecture-notes:page-0020}
\uses{hanlin:thm-mean-value-max}
If $u$ is harmonic on $B_R\setminus\{0\}$ and
\[
u(x)=\begin{cases}o(\log|x|),&n=2,\\o(|x|^{2-n}),&n\ge3,\end{cases}
\quad |x|\to0,
\]
then $u$ extends across $0$ as a $C^2$ harmonic function.
\end{theorem}
\begin{proof}
\sketch Let $v$ be the harmonic extension of the boundary data on $\partial B_R$.
For $w=v-u$, compare $|w|$ on $B_R\setminus B_r$ with the harmonic barriers
$M_rr^{n-2}|x|^{2-n}$ ($n\ge3$) or the logarithmic barrier ($n=2$), where
$M_r=\max_{\partial B_r}|w|$. The assumed little-$o$ decay makes the barrier
tend to zero for each fixed $x\ne0$, so $w=0$ and $u=v$ away from the pole.
\end{proof}

\section{Maximum Principles}

\begin{theorem}\label{hanlin:thm-subharmonic-max}\dcref{Theorem 1.29}\group{harmonic-functions}
\source{han-lin-lecture-notes:page-0020}
\lean{HanLinLectureNotes.Ch01.subharmonic_maximum_principle}
\leanok
If $u\in C^2(B_1)\cap C(\overline{B_1})$ and $\Delta u\ge0$, then
\[
\sup_{B_1}u\le\sup_{\partial B_1}u.
\]
\end{theorem}
\begin{proof}
\sketch For $\varepsilon>0$ set $u_\varepsilon=u+\varepsilon|x|^2$; then
$\Delta u_\varepsilon>0$, so $u_\varepsilon$ has no interior maximum. Compare
its boundary and interior suprema and let $\varepsilon\downarrow0$.
\end{proof}

\begin{corollary}\label{hanlin:corollary-comparison-principle}\dcref{Corollary 1.30}\group{harmonic-functions}
\source{han-lin-lecture-notes:page-0021}
\uses{hanlin:thm-subharmonic-max}
\lean{HanLinLectureNotes.Ch01.comparison_principle}
\leanok
If $u,v\in C^2(B_1)\cap C(\overline{B_1})$ satisfy $\Delta u\ge\Delta v$ in
$B_1$ and $u\le v$ on $\partial B_1$, then $u\le v$ in $B_1$.
\end{corollary}

\begin{remark}\label{hanlin:remark-maxprin-domain}\dcref{Remark 1.31}\group{harmonic-functions}
\source{han-lin-lecture-notes:page-0021}
\uses{hanlin:thm-subharmonic-max, hanlin:corollary-comparison-principle}
Theorem~\ref{hanlin:thm-subharmonic-max} and Corollary~\ref{hanlin:corollary-comparison-principle}
remain valid on every bounded domain, with the corresponding boundary.
\end{remark}

\begin{theorem}\label{hanlin:thm-interior-gradient}\dcref{Theorem 1.32}\group{harmonic-functions}
\source{han-lin-lecture-notes:page-0021}
\uses{hanlin:thm-subharmonic-max, hanlin:lemma-interior-gradient-1}
For harmonic $u$ in $B_1$,
\[
\sup_{B_{1/2}}|Du|\le c\sup_{\partial B_1}|u|,
\qquad
|u(x)-u(y)|\le c|x-y|^\alpha\sup_{\partial B_1}|u|
\]
for $x,y\in B_{1/2}$ and $0\le\alpha\le1$, where $c=c(n)$.
\end{theorem}
\begin{proof}
\sketch $|Du|^2$ is subharmonic since
$\Delta|Du|^2=2\sum_{i,j}(D_{ij}u)^2$. Apply Theorem~\ref{hanlin:thm-subharmonic-max}
to a cutoff multiple of $|Du|^2$ plus a sufficiently large multiple of $u^2$.
The gradient bound gives the Hölder estimate (with a larger constant).
\end{proof}

\begin{lemma}\label{hanlin:lemma-log-gradient}\dcref{Lemma 1.33}\group{harmonic-functions}
\source{han-lin-lecture-notes:page-0022}
If $u>0$ is harmonic in $B_1$, then
\[
\sup_{B_{1/2}}|D\log u|\le c(n).
\]
\end{lemma}
\begin{proof}
\sketch Set $v=\log u$ and $w=|Dv|^2$. Then $\Delta v=-w$ and
$\Delta w+2Dv\cdot Dw=2|D^2v|^2\ge2w^2/n$. Applying the maximum principle to
$\eta^4w$ for a cutoff $\eta$ equal to one on $B_{1/2}$ yields a uniform bound.
\end{proof}

\begin{corollary}\label{hanlin:corollary-harnack-ball}\dcref{Corollary 1.34}\group{harmonic-functions}
\source{han-lin-lecture-notes:page-0023}
\uses{hanlin:lemma-log-gradient}
If $u\ge0$ is harmonic in $B_1$, then
\[
u(x_1)\le c(n)u(x_2)\qquad(x_1,x_2\in B_{1/2}).
\]
\end{corollary}
\begin{proof}
\sketch For $u>0$, integrate the bound in Lemma~\ref{hanlin:lemma-log-gradient}
along the segment joining $x_1$ and $x_2$; the zero case follows by continuity.
\end{proof}

\begin{theorem}\label{hanlin:thm-hopf}\dcref{Theorem 1.35}\group{harmonic-functions}
\source{han-lin-lecture-notes:page-0023,han-lin-lecture-notes:page-0024}
\uses{hanlin:thm-subharmonic-max, hanlin:corollary-harnack-ball}
If $u\in C(\overline{B_1})$ is harmonic in $B_1$ and has a strict maximum at
$x_0\in\partial B_1$, then
\[
\partial_nu(x_0)\ge c(n)\bigl(u(x_0)-u(0)\bigr).
\]
\end{theorem}
\begin{proof}
\sketch On the annulus $B_1\setminus\overline{B_{1/2}}$ use the subharmonic
barrier $v=e^{-\alpha|x|^2}-e^{-\alpha}$ and apply the maximum principle to
$u-u(x_0)+\varepsilon v$. Corollary~\ref{hanlin:corollary-harnack-ball} supplies a
uniform gap on $B_{1/2}$; choosing $\varepsilon$ proportional to that gap and
differentiating the boundary comparison gives the normal derivative bound.
\end{proof}

\begin{lemma}\label{hanlin:lemma-holder-boundary}\dcref{Lemma 1.36}\group{harmonic-functions}
\source{han-lin-lecture-notes:page-0024,han-lin-lecture-notes:page-0025}
\uses{hanlin:corollary-comparison-principle, hanlin:thm-interior-gradient}
If $u\in C(\overline{B_1})$ is harmonic in $B_1$, has boundary data
$\varphi\in C^\alpha(\partial B_1)$ with $0<\alpha<1$, then
$u\in C^{\alpha/2}(\overline{B_1})$ and
\[
\|u\|_{C^{\alpha/2}(\overline{B_1})}\le c(n,\alpha)\|\varphi\|_{C^\alpha(\partial B_1)}.
\]
\end{lemma}
\begin{proof}
\sketch Compare $u-u(x_0)$ with the barrier
$C|x-x_0|^{\alpha/2}$ at each boundary point using Corollary~\ref{hanlin:corollary-comparison-principle}.
For interior points, if their boundary distances are comparable to their
separation use the boundary estimate; otherwise apply the scaled gradient
estimate from Theorem~\ref{hanlin:thm-interior-gradient}. Combining the two cases
gives the global $C^{\alpha/2}$ bound.
\end{proof}

\section{Energy Methods}

Assume $a_{ij}\in C(B_1)$ is uniformly elliptic:
\[
\lambda|\xi|^2\le a_{ij}(x)\xi_i\xi_j\le\Lambda|\xi|^2,
\]
and $u\in C^1(B_1)$ satisfies
\[
\int_{B_1}a_{ij}D_i uD_j\varphi=0\qquad(\varphi\in C_0^1(B_1)).
\]

\begin{lemma}\label{hanlin:lemma-caccioppoli}\dcref{Lemma 1.37}\group{harmonic-functions}
\source{han-lin-lecture-notes:page-0025}
Under these hypotheses, every $\eta\in C_0^1(B_1)$ satisfies
\[
\int_{B_1}\eta^2|Du|^2\le c(\lambda,\Lambda)\int_{B_1}|D\eta|^2u^2.
\]
\end{lemma}
\begin{proof}
\sketch Test the weak equation with $\eta^2u$. Ellipticity bounds the main term,
and Cauchy's inequality absorbs one half of it, leaving the stated estimate.
\end{proof}

\begin{corollary}\label{hanlin:cor-caccioppoli-balls}\dcref{Corollary 1.38}\group{harmonic-functions}
\source{han-lin-lecture-notes:page-0025}
\uses{hanlin:lemma-caccioppoli}
For $0\le r<R\le1$,
\[
\int_{B_r}|Du|^2\le\frac{C(\lambda,\Lambda)}{(R-r)^2}\int_{B_R}u^2.
\]
\end{corollary}
\begin{proof}
\sketch Choose a cutoff equal to one on $B_r$, supported in $B_R$, and with
$|D\eta|\le2/(R-r)$; apply Lemma~\ref{hanlin:lemma-caccioppoli}.
\end{proof}

\begin{corollary}\label{hanlin:cor-energy-decay}\dcref{Corollary 1.39}\group{harmonic-functions}
\source{han-lin-lecture-notes:page-0025,han-lin-lecture-notes:page-0026}
\uses{hanlin:lemma-caccioppoli}
There is $\theta=\theta(n,\lambda,\Lambda)\in(0,1)$ such that, for $0<R\le1$,
\[
\int_{B_{R/2}}u^2\le\theta\int_{B_R}u^2,
\qquad
\int_{B_{R/2}}|Du|^2\le\theta\int_{B_R}|Du|^2.
\]
\end{corollary}
\begin{proof}
\sketch A cutoff in Lemma~\ref{hanlin:lemma-caccioppoli}, followed by Poincare's
inequality on the annulus $B_R\setminus B_{R/2}$, gives the first contraction.
Apply the same argument to $u-a$, choosing $a$ as the annular average, to get
the gradient contraction.
\end{proof}

\begin{remark}\label{hanlin:remark-finite-energy-liouville}\dcref{Remark 1.40}\group{harmonic-functions}
\source{han-lin-lecture-notes:page-0026}
\uses{hanlin:cor-energy-decay}
Iterating Corollary~\ref{hanlin:cor-energy-decay} shows that a harmonic function
on $\mathbb R^n$ with finite $L^2$ norm vanishes, while finite Dirichlet energy
forces it to be constant.
\end{remark}

\begin{remark}\label{hanlin:remark-iterated-energy-decay}\dcref{Remark 1.41}\group{harmonic-functions}
\source{han-lin-lecture-notes:page-0026}
\uses{hanlin:lemma-caccioppoli, hanlin:cor-energy-decay}
For $0<\rho<r\le1$, iteration gives
\[
\int_{B_\rho}u^2\le C(\rho/r)^\mu\int_{B_r}u^2,\qquad
\int_{B_\rho}|Du|^2\le C(\rho/r)^\mu\int_{B_r}|Du|^2,
\]
for some $\mu>0$ depending only on $n,\lambda,\Lambda$ (the gradient exponent
may be chosen in $(n-2,n)$).
\end{remark}

\begin{lemma}\label{hanlin:lemma-constant-coeff-compactness}\dcref{Lemma 1.42}\group{harmonic-functions}
\source{han-lin-lecture-notes:page-0027,han-lin-lecture-notes:page-0028}
\uses{hanlin:cor-caccioppoli-balls}
Let $(a_{ij})$ be a constant positive definite matrix with
\[
\lambda|\xi|^2\le a_{ij}\xi_i\xi_j\le\Lambda|\xi|^2,
\]
and let $u\in C^1(B_1)$ solve the associated weak equation. For
$0<\rho\le r\le1$,
\[
\int_{B_\rho}|u|^2\le c(\rho/r)^n\int_{B_r}|u|^2,
\qquad
\int_{B_\rho}|u-u_\rho|^2\le c(\rho/r)^{n+2}\int_{B_r}|u-\bar u_r|^2,
\]
where $c=c(\lambda,\Lambda)$, and $u_\rho$ and $\bar u_r$ denote the averages
of $u$ on $B_\rho$ and $B_r$, respectively.
\end{lemma}
\begin{proof}
\sketch Scale to $r=1$. Caccioppoli estimates applied successively to the
constant-coefficient equation give local $H^k$ bounds from the $L^2$ norm;
Sobolev embedding yields uniform bounds for $u$ and $Du$ on $B_{1/2}$. Integrate
these bounds over $B_\rho$ and apply the same argument to $u-\bar u_r$ for the
oscillation estimate. Scaling restores the powers of $\rho/r$.
\end{proof}
